\documentclass[11pt]{article}
\usepackage[margin=1in]{geometry}
\usepackage{amsmath,amssymb,amsthm}
\usepackage[hidelinks]{hyperref}
\newtheorem{theorem}{Theorem}
\begin{document}
\title{Ordinary degree-seven binomial SFC reduces to gap six}
\author{}\date{}\maketitle
Let $\mathcal L(X^j)=j!$ and put
\[
 P_j(t)=\sum_{k=0}^j\binom jk(j+1)_{6k}t^k.
\]
Strong Factorial means that every consecutive pair of positive
moments of a nonzero phase has a nonzero entry.

\begin{theorem}[The exact remaining degree-seven family]
\label{thm:gap-six-frontier}
The following are equivalent:
\begin{enumerate}
\item Every ordinary complex univariate binomial of degree at most
seven satisfies Strong Factorial.
\item Every nonzero $X(A+BX^6)$, with arbitrary complex coefficients,
satisfies Strong Factorial.
\item For every terminal integer
\[
 N\ge6,\qquad N\notin\{b7^r:1\le b\le6,\ r\ge0\},
\]
one has $\gcd_{\mathbb Q[t]}(P_{N-1},P_N)=1$.
\end{enumerate}
Every hypothetical common root in the last assertion necessarily
satisfies $v_7(-1/t)=1$ at every embedding into $\mathbb C_7$.
Ordinary supports $\{3,7\},\{4,7\},\{5,7\}$ and every ordinary
binomial of degree at most six are proved at every positive pair.
The remaining gap-six assertion is open.
\end{theorem}
\begin{proof}
Theorem \texttt{thm:two-family-frontier} of the
\href{ORDINARY_DEGREE_SEVEN_TWO_FAMILY_FRONTIER.tex}{two-family classification}
gives the exact residual rows $\{1,7\}$ and $\{5,7\}$.
Corollary \texttt{cor:ordinary-sfc} of the
\href{LOWER_FIVE_QUADRATIC_SECTOR_SFC.tex}{quadratic-sector proof}
completes the second row. This proves the equivalence of the first
and third assertions.

The second assertion is a special case of the first.
Conversely the second assertion supplies every test in the third,
so all three are equivalent. The earlier classification already
covers degree at most six, the other two degree-seven supports,
and all supports through terminal five.
It also covers the displayed omitted index family $N=b7^r$.
The moment normalization is
\[
 \mathcal L((X(A+BX^6))^j)=A^j j!\,P_j(B/A)\qquad(AB\ne0).
\]
A common complex zero is equivalent to a nonconstant rational gcd,
and nonzero monomials have no zero positive moment.

The necessary valuation condition is already part of the two-family
classification, by the
\href{PHYSICAL_PRIME_TOTAL_DEGREE_REDUCTION.tex}{physical prime-total-degree reduction},
\texttt{thm:radius} and the consecutive-moment bridge.
The new sector proof supplies the stated all-index support $\{5,7\}$.
\end{proof}

This is a proved reduction to the remaining ordinary lower-one gap-six
family. It does not prove that family, degree-seven SFC, all-degree
binomial SFC, or unrestricted FC/SFC. The proofs and exact checks have
been reviewed internally; no independent review, formal verification,
priority or novelty claim is made.
\end{document}
